\section{Tổng quan tài liệu}
\label{sec:litreview}

\subsection{Phân loại âm thanh môi trường}

Phân loại âm thanh môi trường (\textit{environmental sound classification}) là
một nhánh của xử lý tín hiệu âm thanh, trong đó mục tiêu là gán nhãn cho một đoạn
tiếng động dựa trên nội dung của nó. So với nhận dạng tiếng nói hay phân loại
nhạc, âm thanh môi trường khó hơn ở chỗ không có cấu trúc lặp rõ ràng như âm vị
hay nhịp điệu, lại thường lẫn nhiều nguồn phát cùng lúc. Bộ dữ liệu UrbanSound8K
được Salamon và cộng sự giới thiệu năm 2014 chính là để chuẩn hóa việc nghiên cứu
mảng này, kèm theo một phân loại học (taxonomy) cho âm thanh đô thị~\cite{salamon2014}.

Điểm chung của các hệ thống phân loại âm thanh là tín hiệu hiếm khi được xử lý ở
dạng sóng thô. Thay vào đó, người ta biến đổi nó sang miền thời gian -- tần số để
làm lộ ra cấu trúc âm học, rồi mới đưa vào bộ phân loại.

\subsection{Hướng tiếp cận truyền thống}

Trước khi học sâu trở nên phổ biến, các hệ thống thường dựa trên đặc trưng thủ
công kết hợp với bộ phân loại cổ điển. Các đặc trưng hay dùng gồm MFCC,
zero-crossing rate, spectral centroid, chroma hay các thống kê năng lượng theo
dải tần -- tất cả đều được thiết kế dựa trên hiểu biết của con người về tín hiệu.
Sau khi trích xuất, đặc trưng được đưa vào SVM, $k$-NN, Random Forest hoặc
Logistic Regression.

Hướng này dễ triển khai và chạy tốt trên dữ liệu nhỏ, nhưng hiệu năng phụ thuộc
nặng vào chất lượng đặc trưng được chọn. Khi dữ liệu đa dạng hơn hoặc ranh giới
giữa các lớp tinh tế hơn, việc thiết kế đặc trưng thủ công nhanh chóng chạm trần.

\subsection{Sự chuyển dịch sang học sâu}

Học sâu thay đổi cục diện bằng cách để mô hình \textit{tự học} biểu diễn đặc
trưng từ dữ liệu thay vì dựa vào đặc trưng do con người định nghĩa. Khi đầu vào
là spectrogram, mạng tích chập có thể khai thác cấu trúc cục bộ trên mặt phẳng
thời gian -- tần số y như cách nó hoạt động trên ảnh: các tầng đầu học những mẫu
đơn giản như biên năng lượng, các tầng sâu hơn tổng hợp chúng thành đặc trưng
mang tính ngữ nghĩa. Piczak~\cite{piczak2015} là một trong những công trình sớm
cho thấy CNN trên log-Mel spectrogram đạt hiệu năng vượt trội so với đặc trưng
thủ công trên chính các bộ dữ liệu âm thanh môi trường.

\subsection{Biểu diễn tín hiệu âm thanh}

\subsubsection{Waveform}

Waveform là dạng biểu diễn gốc theo miền thời gian. Học trực tiếp từ waveform
tránh được mất mát do biến đổi đặc trưng, nhưng đòi hỏi mô hình lớn và quá trình
tối ưu phức tạp -- thường không phù hợp với dữ liệu cỡ vừa và tài nguyên hạn chế.

\begin{figure}[H]
  \centering
  \imgfallback{waveform_siren.png}{0.72\textwidth}
  \caption{Waveform của một đoạn âm thanh thuộc lớp \texttt{siren}.}
  \label{fig:wave}
\end{figure}

\subsubsection{Spectrogram}

Spectrogram mô tả sự phân bố năng lượng của tín hiệu theo thời gian và tần số.
Vì có thể coi như một ảnh hai chiều -- một trục thời gian, một trục tần số -- nó
là biểu diễn rất tự nhiên cho CNN.

\begin{figure}[H]
  \centering
  \imgfallback{spectrogram_siren.png}{0.72\textwidth}
  \caption{Mel spectrogram của cùng đoạn âm thanh \texttt{siren}.}
  \label{fig:spec}
\end{figure}

\subsubsection{Mel spectrogram và log-Mel spectrogram}

Mel spectrogram là spectrogram được chiếu sang thang Mel -- một thang tần số gần
với cảm nhận của tai người hơn thang tuyến tính. Lấy thêm log của năng lượng giúp
nén dải biên độ và làm nổi các vùng năng lượng yếu, cho ra log-Mel spectrogram.
Đây là biểu diễn được dùng phổ biến nhất trong phân loại âm thanh, và cũng là lựa
chọn của đề tài vì nó cân bằng tốt giữa khả năng biểu diễn và độ đơn giản khi
triển khai.

\subsection{CNN cho phân loại âm thanh}

CNN là kiến trúc nền tảng khi xử lý log-Mel spectrogram. Các bộ lọc tích chập học
được những mẫu cục bộ như dải năng lượng ngang, cấu trúc lặp hay biến đổi đột
ngột của tín hiệu. Ưu điểm của một CNN không quá sâu là dễ huấn luyện và ít nhạy
với siêu tham số, nên nó thường được dùng làm \textit{baseline} mạnh -- đề tài
cũng dùng CNN theo đúng vai trò này. Salamon và Bello~\cite{salamon2017} còn cho
thấy rằng khi kết hợp CNN với data augmentation, hiệu năng trên âm thanh đô thị
được cải thiện rõ rệt, đặc biệt ở các lớp vốn ít mẫu hoặc khó.

\subsection{Mạng phần dư ResNet}

ResNet~\cite{he2016} được đề xuất để giải quyết khó khăn khi huấn luyện mạng rất
sâu. Ý tưởng cốt lõi là thay vì để mỗi khối học trực tiếp một ánh xạ mới, mạng
học một \textit{phần dư} cộng vào đầu vào thông qua kết nối tắt (skip connection).
Cách làm này giúp gradient lan truyền thuận lợi và cho phép tăng độ sâu mà không
suy giảm hiệu năng.

Khi đưa ResNet sang miền spectrogram, cần hai điều chỉnh. Đầu vào spectrogram chỉ
có một kênh thay vì ba kênh RGB, nên tầng tích chập đầu phải được sửa lại. Ngoài
ra, các chi tiết cục bộ nhỏ trên spectrogram rất quan trọng, nên phần \textit{stem}
của mạng -- vốn hạ kích thước rất nhanh ở bản gốc cho ảnh -- nên được làm nhẹ hơn
để không xóa mất thông tin sớm. Đề tài dùng ResNet18 vì nó đủ sâu để học đặc
trưng phức tạp hơn CNN baseline nhưng vẫn nhẹ, phù hợp một project cá nhân.

\subsection{Data augmentation trong miền spectrogram}

Augmentation tạo thêm các biến thể hợp lý từ dữ liệu huấn luyện để mô hình học
được đặc trưng bền vững hơn, ít phụ thuộc vào chi tiết ngẫu nhiên của từng mẫu.
Với âm thanh, augmentation có thể thực hiện ở miền waveform hoặc miền spectrogram.

SpecAugment~\cite{park2019} là kỹ thuật tiêu biểu ở miền spectrogram, gồm việc
che một dải tần số (\textit{frequency masking}) hoặc một dải thời gian
(\textit{time masking}) ngay trên biểu diễn hai chiều. Đề tài sử dụng các phép
này cùng với nhiễu Gaussian nhẹ và dịch thời gian. Ngoài ra, MixUp~\cite{zhang2018}
-- trộn tuyến tính hai mẫu và nhãn của chúng -- cũng được áp dụng ở mức batch.
Một lưu ý quan trọng từ tài liệu: nếu augmentation quá mạnh, mô hình có thể bị
huấn luyện trên dữ liệu biến dạng quá mức và giảm hiệu năng -- quan sát này được
kiểm chứng trực tiếp trong thực nghiệm của đề tài.

\subsection{Khoảng trống mà đề tài hướng tới}

Nhiều công trình hiện đại theo đuổi hiệu năng tối đa bằng mô hình rất sâu hoặc mô
hình tiền huấn luyện quy mô lớn -- chẳng hạn PANNs~\cite{kong2020}, được huấn
luyện trước trên AudioSet. Tuy nhiên, trong phạm vi một bài tập lớn, câu hỏi đáng
quan tâm hơn không phải là ``mô hình nào mạnh nhất'' mà là: với cùng một pipeline
minh bạch và tài nguyên vừa phải, các lựa chọn về kiến trúc, augmentation và cách
kết hợp mô hình tác động đến kết quả như thế nào. Đề tài tập trung vào khoảng
trống thực hành đó: so sánh công bằng CNN baseline, ResNet18 và ResNet18 kèm
augmentation trên cùng một dữ liệu và cùng một quy trình, rồi xét xem việc kết
hợp chúng theo hướng ensemble có khai thác được tính bù trừ lỗi hay không.
